\documentclass{gpgc-bs}
\usepackage[utf8]{inputenc}
\usepackage[T1]{fontenc}
\usepackage{geometry}
\usepackage{graphicx}
\usepackage{float}
\usepackage{hyperref}
\usepackage[english,indonesian]{babel}

\title{Optimasi Pemrosesan Template LaTeX Menggunakan Pendekatan Berbasis AI}
\author{Academic Support Agent}
\date{\today}

\begin{document}
\maketitle

\begin{abstract}

\end{abstract}

\chapter{PENDAHULUAN}


\chapter{TINJAUAN PUSTAKA}


% Original content (commented out):
% \documentclass[oneside]{gpgc-bs}
\usepackage{tocloft}
\usepackage{amssymb}
\usepackage{esdiff}

\usepackage{amstext}
\usepackage{algorithmic}
\usepackage{algorithm}
\usepackage{multirow}
\usepackage{rotating}
\usepackage[overload]{textcase}
\usepackage{times}
\usepackage{breqn}
\usepackage{xpatch}
%\usepackage{subfig}
\usepackage{epsfig}
\usepackage{epstopdf}
\usepackage{pdfpages}
\usepackage{rotfloat}
\usepackage{amsmath,amssymb}
\usepackage[all]{xy}
\usepackage{esint}
\usepackage{multicol}
\usepackage{blindtext}
\usepackage{mathptmx}% Times Roman font
\usepackage[scaled=.90]{helvet}% Helvetica, served as a model for arial
\usepackage{xcolor}
\usepackage[T1]{fontenc}
\usepackage{chngcntr}
\usepackage{lipsum} 
\usepackage{rotating}
\usepackage{indentfirst}
\usepackage{sectsty}
\usepackage{titlesec}
\usepackage{mathtools}
\usepackage{amssymb}
\usepackage{amsmath}
\usepackage{notoccite}
\usepackage{comment}
\usepackage{sectsty}
\usepackage{caption}
%\usepackage{numbers}{natbib}
%\usepackage{etoolbox}% 
\usepackage{tocloft}
%\usepackage{subcaption}
\usepackage{subcaption}
\captionsetup{compatibility=false}
\allsectionsfont{\normalsize \bfseries}
\DeclarePairedDelimiter\ceil{\lceil}{\rceil}
\DeclarePairedDelimiter\nint{\lfloor}{\rceil}
\DeclarePairedDelimiter\floor{\lfloor}{\rfloor}
\setcounter{secnumdepth}{10}
\setcounter{tocdepth}{10}
\pagestyle{plain}
\usepackage{titlesec}
\titleformat{\chapter}
{\Large\bfseries\centering\fontsize{14pt}{14pt}\selectfont}
{\MakeUppercase{\chaptertitlename}~\thechapter.}{14pt}{\MakeUppercase}
\titlespacing*{\chapter}{0pt}{0pt}{10pt}
%\setlength{\parindent}{10ex}
 %{\MakeUppercase}
 

\sectionfont{\fontsize{12}{15}\selectfont}
\subsectionfont{\normalfont\itshape\fontsize{12}{15}\selectfont}
\subsubsectionfont{\normalfont\itshape\fontsize{12}{15}\selectfont}
%\subsubsubsectionfont{\itshape\fontsize{12}{15}\selectfont}

\newcommand\countinput[1]{
\input{#1}
\immediate\write18{texcount "#1.tex" -1 -sum > count.txt}
\footnote{FILE \#1 CONTAINS \input{count.txt} WORDS}
}


% Information about the Thesis (All are compulsory, do not delete any line)
%%%%%%%%%%%%%%%%%%%%%%%%%%%%%%%%%%%%%%%%%%%%%%%%%%%%%%%%%%%%%%%%%%%%%%%%%%%%%%%
\university{Hazara University}
\college{Govt. Postgraduate College, Mansehra}
\department{Department of Physics}
\degreeyear{2024}
\degreemonth{April}
\degreename{Physics}
\campuscity{Mansehra}
\author{Fahad Hussain } %keep this space here
\supervisor{Dr. Saeed Ahmad}
%\deanname{}
%\directorname{Dr. Tariq Khan}
\hodname{Prof. Shahid Iqbal}
%\gpcname{Dr. Usman Habib}
\title{Electrochemical evaluation of nanocomposites for efficient energy storage devices}
\rollnumber{Phy-F2}
\internalexaminer{Examiner}
\externalexaminer{Examiner}{Externel Examiner}


% Set of New Symbols Used in this Thesis
%%%%%%%%%%%%%%%%%%%%%%%%%%%%%%%%%%%%%%%%%%%%%%%%%%%%%%%%%%%%%%%%%%%%%%%%%%%%%%%
% If you are writing an oft-repeated symbol again and again, create an entry
% here. 
\newcommand{\fbody}{\ensuremath{\mathbf{f}_{\mathrm{body}}}\xspace}
\newcommand{\functionof}[2]{\ensuremath{#1\left(#2\right)}\xspace}

%%%%%%%%%%%%%%%%%%%%%%%%%%%%%%%%%%%%%%%%%%%%%%%%%%%%%%%%%%%%%%%%%%%%%%%%%%%%%%
% Former document starts below this
\pagenumbering{roman} %for the structural pages



\begin{document}
%\addtocontents{toc}{\vskip -1.2cm}
%\addtocontents{lof}{\vskip -10cm}

\renewcommand{\bibname}{Bibliography}

%\include{Structure/dedication}

%\include{Structure/acknowledgements}

\include{Structure/toc}

%\excludecomment{figure}

%\let\endfigure\relax


\include{Structure/abstract}

\include{Structure/abbreviations}

%\include{Structure/tables-figures}

\mainmatter
%\pagestyle{empty}

\include{Chapters/01-introduction}

\include{Chapters/02-experimental-Setup}

\include{Chapters/03-discussion}

\include{Chapters/04-conclusion}

\bibliographystyle{IEEEtranN}
\addcontentsline{toc}{chapter}{Bibliography}
\bibliography{bibliography} 
\end{document}

%\appendix
%\include{appendix1/appendix1}

\end{document}